% =====================================================================
%  2.2 Likningssystemer — fasit
% =====================================================================

\fk{Del 1 — Innsetting og addisjon}
\fo{1.1} a) Ja. $4 + 3 = 7$ og $4 - 3 = 1$. \ b) Nei. $2 \cdot 3 + 3 = 9$ stemmer, men $3 - 3 = 0 \neq 1$.
\fo{1.2} a) $(3, 5)$ \ b) $(9, 3)$ \ c) $(3, 1)$
\fo{1.3} a) $y$ forsvinner. Vi får $3x = 15$. \ b) $x = 5$ og $y = 4$ \ c) $5 + 4 = 9$ og $2 \cdot 5 - 4 = 6$ stemmer.
\fo{1.4} a) $(4, -1)$ \ b) $(2, -3)$ \ c) $(-2, 5)$
\fo{1.5} a) Innsetting, for $x$ står alene: $(1, 2)$. \ b) Subtraksjon, for $3x$ er lik i begge: $7y = 7$, så $(-3, 1)$.
c) Addisjon, for $6x$ og $-6x$ er motsatte: $6y = 18$, så $(-2, 3)$.
\fo{1.6} Feilen er steget der $4$ flyttes over: $3y = 13 + 4$. Det skal være $3y = 13 - 4 = 9$, så $y = 3$. Riktig svar: $(2, 3)$. Kontroll: $4 + 9 = 13$ og $10 - 12 = -2$.
\fo{1.7} Skolebrød $s$ og kakao $k$: $4s + 2k = 110$ og $3s + 4k = 120$. Svar: skolebrød 20 kr og kakao 15 kr. Kontroll: $80 + 30 = 110$ og $60 + 60 = 120$.
\fo{1.8} Eksempel, flere svar er riktige: $x + y = 2$ og $x - y = 4$. Koeffisientene til $y$ er motsatte, så vi legger sammen: $2x = 6$, $x = 3$, $y = -1$.
Ingen ukjent står alene, så innsetting krever først omskriving.
\fo{1.9} Sett $x = 2$ og $y = 5$ i hver likning. Første: $6 + 5a = 21$, så $a = 3$. Andre: $2b - 5 = 3$, så $b = 4$.
\fo{Sjekk 1} a) $(2, 5)$ \ b) $-4$ \ c) Paret må passe i begge likningene. Et par kan passe i den ene og ikke i den andre.

\fk{Del 2 — Grafisk løsning og antall løsninger}
\fo{2.1} a) $(3, 2)$ \ b) $2 = 3 - 1$ og $2 = -3 + 5$ stemmer. Løsning: $x = 3$ og $y = 2$.
\fo{2.2} a) Ingen løsning (like stigningstall, ulikt konstantledd). \ b) Én løsning (stigningstall 2 og $-1$). \ c) Uendelig mange (samme linje).
\fo{2.3} a) Parallelle \ b) Sammenfallende \ c) De krysser (en løsning for $y$).
\fo{2.4} a) $y = -2x + 5$ og $y = -2x + \frac{7}{2}$. Ingen løsning. \ b) $y = x - 3$ i begge. Uendelig mange løsninger. \ c) Stigningstall $-3$ og $6$. Én løsning: $(1, 1)$.
\fo{2.5} a) $y = 50 + 20x$ og $y = 20 + 30x$. \ b) Etter 3 timer, og det koster 110 kr. \ c) Etter 5 timer: A koster 150 kr, B koster 170 kr. A er billigst, for linje A ligger under linje B etter skjæringspunktet.
\fo{2.6} Feilen er steget $(6x - 3y) - (6x - 3y) \to -6y$. Venstre side er 0, ikke $-6y$. Da står vi igjen med $0 = -6$: ingen løsning. Linjene er parallelle.
Svaret $(3, 1)$ passer i andre likning, men ikke i første: $18 - 3 = 15 \neq 9$.
\fo{2.7} a) $a = 6$ og $b = 2$. Andre likning ganget med 2 gir $2x + 6y = 2b$, og vi trenger $2b = 4$. \ b) $a = 6$ og $b \neq 2$, for eksempel $b = 5$. Da får vi $0 = 2b - 4$, og det er ikke 0. \ c) Alle $a \neq 6$. Gang andre likning med 2 og trekk fra. Da får du $(a - 6)y = 4 - 2b$. Er $a \neq 6$, kan vi dele, og vi får nøyaktig én $y$ og én $x$.
\fo{2.8} a) $x^2 + 2x = 2x + 9$ gir $x^2 = 9$, så $x = -3$ eller $x = 3$. Løsninger: $(-3, 3)$ og $(3, 15)$. Kontroll: $9 - 6 = 3$ og $-6 + 9 = 3$. $9 + 6 = 15$ og $6 + 9 = 15$. \ b) $x^2 = c$. For $c > 0$: to løsninger, $(\pm\sqrt{c},\, c)$. For $c = 0$: én løsning, $(0, 0)$. For $c < 0$: ingen, fordi $x^2 \ge 0$.
\fo{Sjekk 2} a) Ingen løsning (parallelle). \ b) De to likningene er samme linje, så det er uendelig mange løsninger. \ c) Et punkt på linjen passer i likningen. Punktet på begge linjene passer i begge likningene.

\fk{Ekstra utfordring}
\fo{U1} Fra første og tredje: $x + 4 = 6$, så $x = 2$. Så $y = 1$ og $z = 3$. Kontroll: $2 + 1 + 3 = 6$, $2 - 1 = 1$ og $1 + 3 = 4$.
\fo{U2} a) $3v + 2b = 610$ og $2v + 4b = 620$. Gang første med 2 og andre med 3: $6v + 4b = 1220$ og $6v + 12b = 1860$. Trekk fra: $8b = 640$. Barn 80 kr og voksen 150 kr. \ b) Skjæringspunktet er $(150, 80)$.
\fo{U3} Gang første med 6: $3x + 2y = 24$. Sett $x = y + 3$: $5y + 9 = 24$, så $y = 3$ og $x = 6$. Kontroll: $3 + 1 = 4$ og $6 - 3 = 3$.
\fo{U4} a) $aex + bey = ce$ og $bdx + bey = bf$. Trekk fra: $(ae - bd)x = ce - bf$. Del på $ae - bd \neq 0$. \ b) Gang første med $d$ og andre med $a$, og trekk fra: $y = \dfrac{af - cd}{ae - bd}$.
c) Da finnes ikke noen entydig løsning. Stigningstallene $-\frac{a}{b}$ og $-\frac{d}{e}$ er like. Linjene er parallelle eller sammenfallende, så ingen eller uendelig mange løsninger.
\fo{U5} a) Eksempel på kode:
\begin{minipage}{\linewidth}
\begin{kode}
def los(a, b, c, d, e, f):
    D = a*e - b*d
    if D == 0:
        return ("ingen eller uendelig "
                "mange løsninger")
    return ((c*e - b*f)/D, (a*f - c*d)/D)
\end{kode}
\end{minipage}
Test 1 gir $(3, 2)$, skrevet som desimaltall. Test 2 gir teksten. b) Fordel: enkelt og krever ingen algebra. Ulempe: finner bare heltall i et lite område, og bare for systemer med én løsning.
\fo{U6} a) $x^2 = 2x - 1$ gir $(x - 1)^2 = 0$, så $x = 1$ og $y = 1$. Ett skjæringspunkt. \ b) $y = 2x + 3$: $x^2 - 2x + 1 = 4$, $(x - 1)^2 = 4$, $x - 1 = \pm 2$. Punktene er $(-1, 1)$ og $(3, 9)$. Linjen $y = -1$ har ingen, for $x^2 \ge 0$. \ c) Nei. Punktene $(p, p^2)$ og $(q, q^2)$ gir stigningstall $\frac{q^2 - p^2}{q - p} = p + q$. Tre punkt $p < q < r$ på samme linje gir $p + q = q + r$, altså $p = r$. Det er en motsigelse. En loddrett linje $x = k$ treffer parabelen én gang.
